\documentclass[
  12pt,
  a4paper,
  oneside
]{article}

\usepackage{iftex}
\ifPDFTeX
  \usepackage[utf8]{inputenc}
  \usepackage[T1]{fontenc}
  \usepackage{lmodern}
\else
  \usepackage{fontspec}
  \setmainfont{lmroman12-regular.otf}[
    BoldFont=lmroman12-bold.otf,
    ItalicFont=lmroman12-italic.otf
  ]
  \setmonofont{lmmono12-regular.otf}
\fi
\usepackage[ngerman]{babel}
\usepackage{microtype}

\usepackage[
  left=2.5cm,
  right=2.5cm,
  top=2.5cm,
  bottom=2.5cm
]{geometry}
\usepackage{setspace}
\onehalfspacing

\usepackage{array}
\usepackage{booktabs}
\usepackage{tabularx}
\usepackage{amsmath}

\usepackage{csquotes}
\usepackage[
  backend=biber,
  style=authoryear,
  sorting=nyt,
  giveninits=true,
  maxcitenames=2,
  maxbibnames=99
]{biblatex}
\addbibresource{literatur.bib}

\usepackage{graphicx}
\usepackage[
  colorlinks=true,
  linkcolor=black,
  urlcolor=blue,
  citecolor=black
]{hyperref}

\hypersetup{
  pdftitle={Test und Validierung},
  pdfauthor={Daniel Tempel}
}

\setcounter{secnumdepth}{3}
\setcounter{tocdepth}{3}

% Unsichtbarer Platzhalter, damit noch leere Gliederungspunkte im PDF erscheinen.
\newcommand{\contentplaceholder}{\mbox{}\par}

% Erlaubt bei langen technischen Bezeichnern einen Zeilenumbruch nach dem Inhalt.
\newcommand{\inlinecode}[1]{\texttt{#1}\allowbreak}

\begin{document}

\tableofcontents
\clearpage

% Die Nummerierung beginnt bei Kapitel 7.
\setcounter{section}{6}

\section{Test und Validierung}
\label{sec:test-validierung}

Nachdem in den vorangegangenen Kapiteln die Systemarchitektur sowie die
Konzeption und Implementierung des Avisierungsservice und des Webfrontends
beschrieben wurden, wird im Folgenden die entwickelte Lösung systematisch
validiert. Ziel ist es zu überprüfen, ob die zentralen fachlichen und
technischen Anforderungen im implementierten System wie vorgesehen umgesetzt
sind und insbesondere auch relevante Fehler- und Randfälle berücksichtigt
werden.

Die Validierung wird entsprechend der zuvor beschriebenen Systemstruktur
getrennt betrachtet. Abschnitt~\ref{sec:validierung-avisierungsservice-backend}
untersucht den Avisierungsservice anhand automatisierter Tests auf
unterschiedlichen Ebenen, während
Abschnitt~\ref{sec:validierung-frontend} die Validierung des Webfrontends
behandelt. Abschließend werden die erzielten Testergebnisse eingeordnet und die
Grenzen der daraus ableitbaren Aussagen dargestellt.

\subsection{Validierung des Avisierungsservice}
\label{sec:validierung-avisierungsservice-backend}

Die Validierung des Avisierungsservice überprüft ausgewählte fachliche und
technische Eigenschaften der in Kapitel 5 beschriebenen Lösung. Dabei steht
nicht die Prüfung einzelner Klassen im Vordergrund, sondern die Frage, ob
zentrale Anforderungen und Entwurfsentscheidungen unter geeigneten
Prüfbedingungen bestätigt werden können.

Dazu werden unterschiedliche Testebenen eingesetzt. Fachliche Regeln können
weitgehend isoliert geprüft werden, während beispielsweise Mandantentrennung,
Datenbankabfragen oder die Ausführung des BPMN-Prozesses das Zusammenspiel
mehrerer Komponenten erfordern.

\subsubsection{Teststrategie und Testumgebung}
\label{sec:teststrategie-anforderungsbezug-testumgebung}

Die Auswahl des Prüfaufbaus richtet sich danach, welche Eigenschaft untersucht
werden soll. Dabei werden nur diejenigen technischen Komponenten einbezogen,
die für die Aussagekraft des jeweiligen Tests erforderlich sind.

\begin{table}[htbp]
  \centering
  \footnotesize
  \setlength{\tabcolsep}{3pt}
  \renewcommand{\arraystretch}{1.15}
  \caption{Prüfbereiche und zugehöriger Prüfaufbau}
  \label{tab:backend-pruefbereiche}
  \begin{tabularx}{\textwidth}{@{}
    >{\raggedright\arraybackslash}p{0.16\textwidth}
    >{\raggedright\arraybackslash}p{0.25\textwidth}
    >{\raggedright\arraybackslash}p{0.31\textwidth}
    >{\raggedright\arraybackslash}X@{}}
    \toprule
    \textbf{Bereich}
      & \textbf{Repräsentative Prüffälle}
      & \textbf{Erwartetes Verhalten}
      & \textbf{Prüfaufbau} \\
    \midrule
    Domänen- und Anwendungslogik
      & abgelaufene Anfrage; identischer Avisierungsauftrag; unzulässiger
        Zustandswechsel
      & Fachregeln werden eingehalten und ungültige Zustandsänderungen verhindert
      & Tests der Domänen- und Anwendungslogik mit kontrollierten Abhängigkeiten \\
    REST-Verträge
      & fehlende Pflichtangaben; ungültige Eingaben; fachlich zurückgewiesene Anfragen
      & definierte HTTP-Statuscodes und Fehlerantworten
      & MVC-Tests mit \texttt{MockMvc} \\
    Zugriffsschutz und Mandantentrennung
      & ungültiger API-Key; Zugriff auf fremde Daten; geschützter Zugriff auf das
        Avisierungsdashboard
      & unzulässige Zugriffe werden verhindert und Mandantengrenzen eingehalten
      & Spring-Boot-Integrationstests mit Security und Persistenz \\
    Persistenz und Abfragen
      & aktuelle und historische Anfragen; Filterung nach Tour, Status und Zeitraum
      & Daten werden korrekt gespeichert, ausgewählt, gruppiert und gefiltert
      & PostgreSQL-, JPA- und QueryDSL-Integrationstests \\
    Workflow und Nebenläufigkeit
      & Versand; Fristablauf; Wiederholungen; parallele Aktionen; veraltete Anfrage
      & Prozesspfade werden korrekt ausgeführt und Zustände bleiben konsistent
      & Camunda- und Concurrency-Integrationstests \\
    Ausgehende Kommunikation
      & Benachrichtigungsversand; DISPO-Callback; Fehler externer Abhängigkeiten
      & ausgehende Aufrufe und Fehlerverarbeitung entsprechen der vorgesehenen
        Logik des Avisierungsservice
      & Adapter- und Integrationstests mit kontrollierten Antworten \\
    \bottomrule
  \end{tabularx}
\end{table}

Für isolierte Prüfungen werden nicht relevante Abhängigkeiten durch Test
Doubles ersetzt. REST-Verträge werden mit \texttt{MockMvc} geprüft.
Sicherheitsrelevante Szenarien verwenden einen erweiterten
Spring-Boot-Testkontext mit aktivierter Security-Konfiguration und Persistenz.

Datenbankbezogene Prüfungen werden mit PostgreSQL über Testcontainers
durchgeführt. Das Schema wird mit den Flyway-Migrationen aufgebaut, sodass
JPA-Mappings, Persistenzverhalten und QueryDSL-Abfragen gegen dieselbe
Datenbanktechnologie wie im Avisierungsservice geprüft werden.

Für die Workflow-Validierung wird die ausführbare BPMN-Prozessdefinition in den
Testkontext eingebunden. Die automatische Camunda-Jobausführung ist deaktiviert,
sodass Jobs und zeitabhängige Abläufe gezielt und reproduzierbar ausgeführt
werden können. Externe Benachrichtigungs- und DISPO-Komponenten werden dabei
durch Test Doubles ersetzt, da der interne Workflow und nicht die Verfügbarkeit
externer Systeme geprüft wird.

Die automatisierten Tests des Avisierungsservice werden im Rahmen des
Maven-Builds ausgeführt. Zusätzlich sind sie in eine GitHub-Actions-Pipeline
eingebunden. Bei Änderungen unter \texttt{backend/**} oder an der zugehörigen
CI-Workflow-Datei wird für Pushes und Pull Requests gegen den Hauptbranch ein
Maven-\texttt{verify}-Lauf unter Java~17 ausgeführt. Dadurch wird die
Testausführung unabhängig von der lokalen Entwicklungsumgebung reproduzierbar
in den Integrationsprozess eingebunden. Die folgenden Abschnitte betrachten die
einzelnen Prüfbereiche anhand ausgewählter Szenarien genauer.

\subsubsection{Domänen- und Anwendungslogik}
\label{sec:test-domaenen-anwendungslogik}

Die Tests der Domänen- und Anwendungslogik prüfen ausgewählte fachliche
Invarianten und Entscheidungsregeln des Avisierungsservice. Im Mittelpunkt
stehen zeitliche Grenzbedingungen, Zustandsänderungen sowie die
Wiederverwendung, Aktualisierung oder Ersetzung bestehender
Rückbestätigungsanfragen.

Zeitabhängige Regeln werden gezielt an ihren Grenzen geprüft. Ein Lieferfenster
muss bei der Übernahme noch in der Zukunft liegen, und eine Kundenreaktion ist
nur vor Ablauf der Antwortfrist zulässig. Betrachtet werden daher Zeitpunkte
unmittelbar vor sowie beim Erreichen der jeweiligen Grenze.

Bei Rückbestätigungsanfragen werden insbesondere unzulässige Zustandsänderungen
geprüft. Eine bereits als \texttt{SENT} markierte Anfrage darf weder erneut in
diesen Zustand überführt noch nachträglich als Zustellfehler gekennzeichnet
werden. Änderungen der Auftragsdaten verändern außerdem nicht automatisch den
fachlichen Bestätigungsstatus.

Für erneut übermittelte Avisierungsaufträge wird geprüft, ob bestehende Anfragen
weiterverwendet, aktualisiert oder ersetzt werden. Eine noch nicht versendete
Anfrage kann bei geänderten Daten aktualisiert werden. Bei einer aktiven
\texttt{SENT}-Anfrage wird dagegen die bestehende Anfrage deaktiviert, ihre
Ablösung gemeldet und eine Neuerzeugung mit den aktualisierten Daten veranlasst.
Unveränderte laufende oder bereits bestätigte beziehungsweise abgelehnte
Aufträge werden wiederverwendet, während nach Versandfehler oder
\texttt{NO\_RESPONSE} auch bei unveränderten Daten eine neue Anfrage entstehen
kann.

Bei Kundenreaktionen werden neben Bestätigung und Ablehnung auch ungültige Fälle
betrachtet. Antworten auf inaktive, abgelaufene, veraltete oder bereits
beantwortete Anfragen werden zurückgewiesen. Für die Duplikatkontrolle wird eine
bereits vorhandene Kundenreaktion simuliert. Dabei bleiben Auftragsstatus und
Aktivitätszustand unverändert; weder die Speichermethode des Repositorys noch
der Workflow-Port werden erneut aufgerufen.

Der Fristablauf darf den Auftrag nur dann auf \texttt{NO\_RESPONSE} setzen, wenn
die Anfrage noch aktiv und die Antwortfrist erreicht ist. Bei einem vorzeitigen
Timeout oder einer bereits inaktiven Anfrage erfolgt keine Statusänderung.

Auch die manuelle erneute Avisierung wird anhand ihrer fachlichen
Voraussetzungen geprüft. Sie ist nach Versandfehler oder \texttt{NO\_RESPONSE}
zulässig, nicht jedoch bei noch ausstehenden, aktiven oder bereits beantworteten
Anfragen. Zusätzlich wird geprüft, dass ein unzulässiger Zustand oder ein
bereits begonnenes Lieferfenster vor nachgelagerten Bedingungen erkannt wird.
Wird die erneute Avisierung aufgrund eines unzulässigen Zustands zurückgewiesen,
bleiben der bestehende Auftragsstatus und die vorhandenen
Rückbestätigungsanfragen unverändert.

Der folgende Abschnitt betrachtet die Eingabevalidierung sowie die Abbildung
fachlicher Ergebnisse und Fehler auf den REST-Vertrag.

\subsubsection{REST-Verträge, Eingabevalidierung und Fehlerantworten}
\label{sec:test-rest-vertraege-fehlerantworten}

Die Prüfung der REST-Schnittstellen untersucht, ob HTTP-Anfragen korrekt
verarbeitet und Ergebnisse der Anwendungsschicht konsistent in HTTP-Antworten
überführt werden. Dafür werden sowohl isolierte MVC-Tests mit ersetzten
Anwendungsfällen als auch Spring-Boot-Integrationstests eingesetzt. Während
MVC-Tests vor allem Parameterbindung, Validierung und Controller-Mapping prüfen,
beziehen Integrationstests zusätzlich ausgewählte Teile der Anwendungslogik und
Persistenz ein.

Für gültige Abfragen des Avisierungsdashboards wird geprüft, ob Statusfilter,
Tournummern und Datumsgrenzen korrekt in das Abfrageobjekt übernommen werden.
Ungültige Parameter wie ein unbekannter Statuswert oder ein nicht ISO-konformes
Datum führen zu \texttt{400 Bad Request}, ohne dass der Anwendungsfall
aufgerufen wird.

Zusätzlich wird die DTO-Validierung betrachtet. Fehlt beispielsweise in einem
syntaktisch gültigen JSON-Request das Pflichtfeld \texttt{responseType}, wird
dieser mit \texttt{400 Bad Request} und \texttt{VALIDATION\_ERROR}
zurückgewiesen. Fachlich ungültige Daten werden dagegen durch die
Anwendungslogik geprüft. Ein Lieferfenster in der Vergangenheit führt ebenfalls
zu \texttt{400 Bad Request}; anschließend sind weder Auftrag noch
Rückbestätigungsanfrage gespeichert.

Davon zu unterscheiden ist die Abbildung bereits festgestellter fachlicher
Fehler. Im MVC-Test für einen ungültigen Zeitraum einer Abfrage des
Avisierungsdashboards wird der ersetzte Anwendungsfall so konfiguriert, dass er
eine \texttt{InvalidFilterException} auslöst. Geprüft wird deren Abbildung auf
\texttt{400 Bad Request} mit dem Fehlercode \texttt{VALIDATION\_ERROR}.

Auch erfolgreiche sowie fachlich zurückgewiesene Anfragen werden auf HTTP-Ebene
geprüft. Ein übernommener DISPO-Auftrag mit Status \texttt{OPEN} liefert
\texttt{202 Accepted}. Der Status bestätigt lediglich die Annahme zur weiteren
Verarbeitung, nicht den Abschluss des Avisierungsprozesses. Eine gültige
Kundenreaktion führt zu \texttt{204 No Content}, eine abgelaufene Anfrage zu
\texttt{410 Gone} und eine inaktive beziehungsweise bereits verarbeitete
Anfrage zu \texttt{409 Conflict}.

Für die Detailansicht des Avisierungsdashboards wird außerdem geprüft, ob
Auftragsdaten, aktuelle Rückbestätigungsanfrage, Historie und Kundenreaktion in
der vorgesehenen JSON-Struktur ausgegeben werden. Fehlerantworten enthalten einen
maschinenlesbaren Fehlercode, eine Beschreibung, HTTP-Status, Pfad und
Zeitstempel. Für einen nicht vorhandenen Auftrag werden alle fünf Felder
geprüft; weitere Szenarien kontrollieren insbesondere die Zuordnung von Status
und Fehlercode. Die gemeinsame Struktur wird durch den
\texttt{GlobalExceptionHandler} erzeugt.

Auch die E-Mail-Einstellungen prüfen weitere Vertragseigenschaften. Fehlt eine
SMTP-Konfiguration, wird \texttt{200 OK} mit \texttt{configured=false}
zurückgegeben. Passwort und verschlüsselte Repräsentation dürfen nicht
Bestandteil der Antwort sein. Werden die ersetzten Anwendungsfälle für
Verbindungstest oder Testnachricht mit den vorgesehenen Ausnahmen konfiguriert,
wird deren Abbildung auf \texttt{502 Bad Gateway} geprüft.

In den MVC-Slice-Tests sind Anwendungsfälle ersetzt und Security-Filter
deaktiviert. Integrationstests beziehen ausgewählte Anwendungs- und
Persistenzkomponenten ein, ersetzen jedoch ebenfalls bestimmte Abhängigkeiten
durch Test Doubles. Abschnitt~\ref{sec:test-zugriffsschutz-mandantenisolation}
betrachtet daher ergänzend Zugriffsschutz, Sitzungsverwaltung und
Mandantenisolation. Die übergreifenden Aussagegrenzen werden in
Abschnitt~\ref{sec:test-ergebnisse-aussagegrenzen} zusammengefasst.

\subsubsection{Zugriffsschutz, Sitzungsverwaltung und Mandantenisolation}
\label{sec:test-zugriffsschutz-mandantenisolation}

Während Abschnitt~\ref{sec:test-rest-vertraege-fehlerantworten} die Verarbeitung
von HTTP-Anfragen und deren Abbildung auf den REST-Vertrag untersucht,
betrachtet dieser Abschnitt die Durchsetzung von Zugriffsregeln und
Mandantengrenzen. Hierfür werden Integrationstests mit aktiver
Sicherheitskonfiguration und PostgreSQL eingesetzt.

Für die DISPO-Schnittstelle wird geprüft, dass fehlende oder ungültige
Zugangsdaten zurückgewiesen und gültige Zugangsdaten dem richtigen Unternehmen
zugeordnet werden. Gleichzeitig wird sichergestellt, dass die für die
System-zu-System-Kommunikation verwendeten Zugangsdaten keinen direkten Zugriff
auf geschützte Funktionen des Avisierungsdashboards ermöglichen.

Der Zugriff auf das Avisierungsdashboard wird über einen zeitlich begrenzten
Zugangscode in eine authentifizierte Sitzung überführt. Die Tests prüfen, dass
ein gültiger Code erfolgreich verwendet werden kann, während eine erneute
Verwendung oder ein bereits abgelaufener Code zurückgewiesen wird.

Für geschützte Operationen des Avisierungsdashboards wird zusätzlich der
CSRF-Schutz überprüft. Ändernde Anfragen ohne gültigen CSRF-Nachweis werden
abgewiesen, bevor die zugehörige Anwendungslogik ausgeführt wird. Bei gültiger
Authentifizierung und gültigem CSRF-Nachweis wird die vorgesehene Operation
zugelassen.

Der Kundenzugriff erfolgt unabhängig von den Authentifizierungsmechanismen für
DISPO und das Avisierungsdashboard über einen individuellen Zugriffstoken. Die
Tests prüfen, dass gültige Tokens die zugehörigen Bestätigungsdaten bereitstellen,
unbekannte oder durch eine neuere Anfrage abgelöste Tokens dagegen keinen Zugriff
mehr ermöglichen.

Ein weiterer Schwerpunkt liegt auf der Mandantenisolation. Hierzu werden Daten
mehrerer Unternehmen parallel angelegt und anschließend geprüft, dass Lese- und
Änderungsoperationen jeweils auf die Daten des authentifizierten Unternehmens
beschränkt bleiben. Fremde Aufträge können nicht verändert werden. Dieselbe
Trennung wird zusätzlich für unternehmensspezifische E-Mail-Konfigurationen
überprüft.

Neben der Zugriffskontrolle werden ausgewählte kryptographische
Schutzmechanismen und die Token-Erzeugung isoliert geprüft. Für die
Verschlüsselung sensibler SMTP-Konfigurationsdaten werden erfolgreiche Ver- und
Entschlüsselung, die Verwendung unterschiedlicher Initialisierungsvektoren sowie
die Bindung des Chiffrats an den jeweiligen Unternehmenskontext überprüft. Die
Token-Erzeugung wird hinsichtlich URL-sicherer Darstellung und der erwarteten
Länge des zugrunde liegenden Zufallswerts geprüft. Zusätzlich wird kontrolliert,
dass zwei aufeinanderfolgend erzeugte Tokens unterschiedliche Werte besitzen.

Die Tests bestätigen damit die Trennung der verschiedenen Zugriffsmodelle sowie
die Mandantenisolation für ausgewählte Lese- und Schreiboperationen. Die
Aussagegrenzen der Sicherheitsprüfungen werden gemeinsam mit den übrigen
Testgrenzen in Abschnitt~\ref{sec:test-ergebnisse-aussagegrenzen}
zusammengefasst. Der folgende Abschnitt betrachtet die Persistenz und die
leseseitigen Abfragen.

\subsubsection{Persistenz- und Abfragemodell}
\label{sec:test-persistenz-abfragemodell}

Die Persistenz- und Abfragetests prüfen, ob Avisierungsdaten entsprechend dem
Datenmodell gespeichert und für Tourübersicht und Detailansicht korrekt
rekonstruiert werden. Als Testdatenbank wird PostgreSQL über Testcontainers
verwendet. Das Schema wird durch Flyway aufgebaut und anschließend von
Hibernate gegen die JPA-Abbildungen validiert. Vor den eigentlichen Abfragen
werden ausstehende Änderungen mit \texttt{flush} in die Datenbank geschrieben
und der Persistence Context geleert. Dadurch basieren die Ergebnisse auf
tatsächlich persistierten Daten.

Eine zentrale Regel des Abfragemodells besteht darin, für einen Auftrag stets
die Rückbestätigungsanfrage mit der höchsten internen ID als aktuelle Anfrage zu
verwenden. Die in
Abschnitt~\ref{sec:test-zugriffsschutz-mandantenisolation} betrachtete
Zugriffsbeschränkung wird auf Persistenzebene dadurch unterstützt, dass
ausschließlich die jeweils neueste Anfrage eines Auftrags aufgelöst wird. Diese
Auswahl erfolgt jeweils innerhalb des zugehörigen Auftrags. Auch Tourübersicht
und Detailansicht verwenden diese Regel, sodass Lieferdatum, Lieferfenster und
Kommunikationskanal aus der jeweils aktuellen Anfrage stammen.

Für die Tourübersicht werden Filterung und Suche separat betrachtet.
Datumsgrenzen werden einschließlich ihrer Randwerte ausgewertet.
Unterschiedliche Filterkriterien werden miteinander verknüpft, während mehrere
Werte innerhalb eines Kriteriums alternativ gelten. Zusätzlich wird zwischen
Suchtreffern auf Tour- und Auftragsebene unterschieden. Ein Treffer in
Tournummer oder Kennzeichen lässt die Tour grundsätzlich im Ergebnis, während
weitere Filter weiterhin auf die enthaltenen Aufträge wirken. Ein Treffer
ausschließlich in auftragsbezogenen Feldern führt dagegen dazu, dass nur die
passenden Aufträge innerhalb der Tour angezeigt werden. Die Suche ist
unabhängig von Groß- und Kleinschreibung.

Für Gruppierung und Paginierung gilt die fachliche Vorgabe, dass eine Tournummer
nicht für unterschiedliche Liefertage oder Fahrzeuge wiederverwendet wird.
Aufträge derselben Tour werden daher gemeinsam dargestellt. Die Touren werden
nach Lieferdatum und Tournummer sortiert, die enthaltenen Aufträge nach Beginn
und Ende des Lieferfensters sowie nach externer Auftrags-ID. Die Paginierung
erfolgt auf Tour-Ebene und trennt damit keine Aufträge derselben Tour. In einem
Testszenario werden drei Touren bei einer Seitengröße von zwei auf zwei Seiten
verteilt; sämtliche Aufträge der ersten Tour bleiben dabei gemeinsam erhalten.

Für die Detailansicht wird geprüft, ob die aktuelle Anfrage von ihrer Historie
korrekt getrennt wird. Die neueste Anfrage erscheint als
\texttt{currentRequest}, ältere Anfragen werden in absteigender Reihenfolge
ihrer ID als \texttt{previousRequests} ausgegeben. Diese Zuordnung hängt nicht
davon ab, ob die aktuelle Anfrage bereits versendet wurde oder noch den
Versandstatus \texttt{PENDING} beziehungsweise \texttt{FAILED} besitzt.
Vorhandene Kundenreaktionen werden mit Typ, Kommentar und Eingangszeitpunkt
übernommen. Für Aufträge ohne Rückbestätigungsanfragen wird ein Ergebnis ohne
\texttt{currentRequest} und mit leerer Historie geliefert; bei einem unbekannten
Auftrag wird kein Detailergebnis zurückgegeben. Die Mandantenisolation wurde
bereits in Abschnitt~\ref{sec:test-zugriffsschutz-mandantenisolation}
betrachtet.

Die Tests bestätigen damit die wesentlichen Persistenz- und Abfrageregeln unter
Verwendung derselben Datenbanktechnologie wie im Avisierungsservice.

\subsubsection{Workflow-, Wiederholungs- und Nebenläufigkeitsszenarien}
\label{sec:test-workflow-wiederholung-nebenlaeufigkeit}

Die Workflow-Tests prüfen ausgewählte Pfade der ausführbaren
BPMN-Prozessdefinition sowie das Verhalten bei Wiederholungen und
konkurrierenden Operationen. Die automatische Job-Ausführung ist dabei
deaktiviert, sodass einzelne Jobs und Timer gezielt ausgeführt und die daraus
entstehenden Prozess- und Datenbankzustände überprüft werden können. Für
zeitabhängige Szenarien werden Anwendungs- und Prozesszeit kontrolliert gesetzt.

Für den regulären Prozessablauf wird zunächst geprüft, dass nach dem
Prozessstart ein asynchroner Versand-Job vorliegt und noch keine Versandaktion
erfolgt ist. Nach dessen erfolgreicher Ausführung muss die
Rückbestätigungsanfrage den Versandstatus \texttt{SENT} und der Auftrag den
Bestätigungsstatus \texttt{SENT} besitzen. Zusätzlich werden die vorgesehenen
Wartezustände des Prozesses sowie der Fälligkeitszeitpunkt des Antworttimers
überprüft.

Weitere Tests betrachten die alternativen Prozesspfade. Eine Kundenreaktion
muss einen Callback-Job mit den erwarteten Prozessdaten erzeugen und der Prozess
nach erfolgreicher Verarbeitung enden. Beim Erreichen der Antwortfrist werden
Anfrage und Auftrag auf den vorgesehenen Zustand \texttt{NO\_RESPONSE}
überführt und ebenfalls ein Callback vorbereitet. Wird eine Anfrage durch eine
neuere ersetzt, muss die bisherige Prozessinstanz dagegen ohne Callback beendet
werden.

Die Wiederholungslogik des Versands wird mit kontrolliert erzeugten Fehlern
geprüft. Bei einem wiederholbaren Fehler werden insgesamt bis zu drei
Versandversuche durchgeführt. Die Tests überprüfen dabei sowohl die Anzahl der
Versuche als auch die konfigurierten Wartezeiten von einer beziehungsweise fünf
Minuten. Bleiben alle Versuche erfolglos, wird die Anfrage auf \texttt{FAILED}
gesetzt und der Prozess ohne DISPO-Callback beendet. Für einen als dauerhaft
behandelten Fehler wird geprüft, dass keine weiteren fachlich modellierten
Wiederholungsversuche stattfinden.

Davon getrennt wird das Verhalten bei einem unerwarteten technischen Fehler
geprüft. In diesem Fall darf kein fachlicher Versandversuch fortgeschrieben
werden. Stattdessen wird der technische Retry-Zähler des Camunda-Jobs reduziert,
während die Anfrage im vorherigen Zustand verbleibt und kein fachlicher
Wiederholungstimer entsteht. Dadurch wird geprüft, dass fachliche und technische
Wiederholungsmechanismen voneinander getrennt bleiben.

Auch Fehler beim DISPO-Callback werden untersucht. Nach Ausschöpfung der
technischen Job-Retries muss ein \texttt{failedJob}-Incident vorliegen, während
die Prozessinstanz bestehen bleibt. Ein zusätzlicher Test prüft den Fall, dass
die betreffende Rückbestätigungsanfrage zwischenzeitlich durch eine neuere
ersetzt wurde. Der veraltete Callback darf dann nicht mehr ausgeführt werden;
die alte Prozessinstanz und der zugehörige Incident werden beendet.

Die Nebenläufigkeitstests prüfen mehrere konkurrierende Zugriffe auf denselben
fachlichen Zustand. Bei zwei parallelen Versandaufrufen wird kontrolliert, dass
die externe Versandaktion nur einmal erfolgt und anschließend ein konsistenter
Versandstatus gespeichert ist. Der Test bezieht sich dabei auf parallele
Aufrufe der Anwendungslogik und nicht auf mehrere gleichzeitig arbeitende
Camunda Job Executor.

Für gleichzeitig eintreffende Kundenreaktionen \texttt{CONFIRM} und
\texttt{REJECT} wird geprüft, dass genau eine Reaktion gespeichert und nur eine
der beiden Operationen erfolgreich abgeschlossen wird. Der resultierende
Auftragsstatus muss mit der tatsächlich gespeicherten Reaktion übereinstimmen;
welche der konkurrierenden Reaktionen gewinnt, ist nicht festgelegt.

Ein weiterer Nebenläufigkeitstest betrachtet die gleichzeitige
Kundenbestätigung und Ablösung einer bestehenden Anfrage. Dabei wird geprüft,
dass die beiden Operationen nicht zu einem widersprüchlichen Zwischenzustand
führen. Nach Abschluss beider Transaktionen ist genau eine Kundenbestätigung
gespeichert, die bisherige Anfrage inaktiv und der Auftrag für die neu erzeugte
Avisierung wieder \texttt{OPEN}.

\subsubsection{Ausgehende Kommunikation und Adapterverträge}
\label{sec:test-kommunikation-adaptervertraege}

Die Tests der ausgehenden Kommunikation prüfen, ob der Avisierungsservice den
für die Rückbestätigungsanfrage festgelegten Kommunikationskanal korrekt auf den
zugehörigen Adapter abbildet, die erforderlichen Daten übergibt und ausgewählte
Fehler an den Systemgrenzen korrekt verarbeitet. Im Unterschied zu
Abschnitt~\ref{sec:test-workflow-wiederholung-nebenlaeufigkeit} steht dabei die
backendseitige Schnittstelle zu externen Systemen im Mittelpunkt. Externe
Sendeoperationen werden ersetzt, HTTP-Antworten durch das Testsetup vorgegeben.

Für E-Mail, SMS und WhatsApp wird geprüft, dass der jeweils vorgesehene Sender
aufgerufen und mit den relevanten Auftrags- und Anfragedaten versorgt wird. Dazu
wird zunächst ein Avisierungsauftrag angelegt und anschließend der Versand
ausgelöst. Durch die ersetzten Sendeoperationen lassen sich Kanalzuordnung und
Datenübergabe gezielt beobachten.

Beim E-Mail-Versand wird geprüft, dass die mandantenspezifischen
SMTP-Einstellungen anhand der Unternehmens-ID angefordert und für die
Senderkonfiguration verwendet werden. Separate Tests untersuchen ausgewählte
Eigenschaften von \texttt{STARTTLS} und Implicit TLS sowie das Verhalten bei
deaktivierter SMTP-Authentifizierung. Ohne Authentifizierung werden keine
Zugangsdaten verwendet und keine Passwortentschlüsselung durchgeführt.

Für die initiale Rückbestätigungsanfrage sowie die Eingangsbestätigungen nach
Bestätigung oder Ablehnung werden Adressierung, Nachrichtenvariante und
ausgewählte Template-Inhalte geprüft. Dabei werden Absender, Empfänger und
Betreff sowie die Übergabe der Bestätigungs-URL kontrolliert. Im initialen
Template werden Lieferdatum und Lieferfenster geprüft; alle drei
Nachrichtenvarianten enthalten zudem die vorgesehene Referenz auf das
eingebettete Logo.

Fehlen erforderliche SMTP-Einstellungen, wird der Versand abgebrochen, bevor
eine Nachricht versendet wird. Ein simulierter Fehler beim eigentlichen
E-Mail-Versand wird in einen für die Anwendung definierten Versandfehler
überführt, wobei die ursprüngliche technische Ursache erhalten bleibt. Damit
wird ein ausgewählter Fehlerpfad geprüft, nicht die vollständige
E-Mail-Fehlerbehandlung.

Für den DISPO-Callback wird ein HTTP-POST an die übergebene Zieladresse mit
Auftrags-ID, Bestätigungsstatus und Kundenkommentar geprüft. Eine erfolgreiche
Antwort wird korrekt verarbeitet. Ein simulierter Serverfehler führt dagegen
zu einem für die Anwendung definierten Callback-Fehler, wobei der zugehörige
Auftrags- und Statuskontext erhalten bleibt.

Die Ergebnisse dieser Adapterprüfungen werden gemeinsam mit den übrigen
Testbereichen im folgenden Abschnitt eingeordnet.

\subsubsection{Ergebnisse und Aussagegrenzen der Tests des Avisierungsservice}
\label{sec:test-ergebnisse-aussagegrenzen}

Die Ergebnisse stützen die korrekte Umsetzung zentraler fachlicher Regeln und
technischer Schutzmechanismen für die untersuchten Szenarien. Unzulässige
Zustandsänderungen werden zurückgewiesen, zeitliche Zulässigkeitsgrenzen
berücksichtigt und gültige Kundenreaktionen, Fristabläufe sowie erneute
Avisierungen in die vorgesehenen Zustände überführt. REST-Schnittstellen,
Zugriffsregeln und Mandantengrenzen verhalten sich in den geprüften Fällen
entsprechend den definierten Vorgaben.

Auf Integrationsebene liefern die PostgreSQL-Abfragen die erwarteten Filter-,
Gruppierungs-, Historien- und Paginierungsergebnisse. Die geprüften BPMN-Pfade
verarbeiten Kundenreaktionen, Fristabläufe, Versandwiederholungen und abgelöste
Anfragen wie vorgesehen. In den untersuchten Konkurrenzsituationen bleiben die
geprüften Konsistenzbedingungen erhalten. Die ausgehenden Adapter erzeugen die
erwarteten Nachrichten und HTTP-Aufrufe und bilden ausgewählte technische Fehler
auf die vorgesehene Anwendungslogik des Avisierungsservice ab.

Für den dokumentierten Softwarestand wurde der vollständige
Maven-\texttt{verify}-Lauf erfolgreich abgeschlossen.

Die Aussagekraft ist auf die beschriebenen Prüfbedingungen begrenzt. Externe
Systeme wie DISPO, SMTP-Provider oder Twilio werden teilweise durch Test Doubles
oder simulierte Antworten ersetzt. Geprüft werden damit die backendseitigen
Adapterverträge, nicht die tatsächliche Verfügbarkeit oder das Verhalten
produktiver Fremdsysteme. Reales Browserverhalten, Cookie-Eigenschaften und der
konfigurierte Inaktivitätstimeout werden ebenfalls nicht abgedeckt.

Camunda-Jobs und Timer werden bei deaktivierter automatischer Jobausführung
gezielt ausgeführt. Daher werden kontrollierte Prozessabläufe geprüft, nicht das
vollständige Betriebsverhalten der automatischen Jobausführung. Ebenfalls nicht
untersucht werden Last- und Performanzverhalten, die Migration bestehender
Produktivdatenbanken, mehrere parallel betriebene Instanzen des
Avisierungsservice oder beliebige Ausfallszenarien. Die Nebenläufigkeitstests
decken nur ausgewählte Konkurrenzsituationen ab.

Eine weitere Grenze betrifft externe Seiteneffekte. Da Datenbankänderungen und
externe Kommunikationsaufrufe nicht atomar gekoppelt sind, kann eine erneute
Ausführung zu wiederholten externen Effekten führen. Die implementierten
Zustandsprüfungen reduzieren dieses Risiko, gewährleisten jedoch keine globale
Exactly-once-Ausführung.

Insgesamt bestätigen die Tests, dass die zentralen fachlichen Regeln,
Sicherheitsmechanismen sowie Persistenz-, Workflow- und Integrationsfunktionen
des Avisierungsservice unter den definierten Prüfbedingungen wie vorgesehen
zusammenwirken. Damit stützen die Testergebnisse die technische Tragfähigkeit
des in Kapitel 5 beschriebenen Avisierungsservice.

\subsection{Validierung des Webfrontends}
\label{sec:validierung-frontend}

% Zuständigkeit: anderes Teammitglied.
% TODO: Inhalt durch das zuständige Teammitglied ergänzen.
\contentplaceholder

\end{document}
